\chapter{Modéliser : le laboratoire de modèles métier}
\label{ch:metier}

\questionchapitre{Avant d'invoquer l'apprentissage profond, que sait déjà faire la modélisation
hydrogéologique établie, et comment l'outiller pour qu'elle serve de référence ?}

\section{Pourquoi commencer par le modèle métier}

Il aurait été possible d'aller directement à l'apprentissage profond. Ce choix aurait été
méthodologiquement fautif, pour trois raisons.

Sans référence, une performance ne signifie rien. Annoncer qu'un réseau de neurones prévoit un
niveau de nappe avec une certaine erreur n'a de sens que comparé à ce qu'obtient la méthode
déjà employée par les hydrogéologues. Une architecture profonde qui n'égale pas un modèle à
quelques paramètres ne mérite pas d'être déployée.

Ensuite, le modèle métier est interprétable par construction, et fournit donc un point d'ancrage
au travail d'explicabilité. On sait ce que ses paramètres représentent, et l'on peut confronter
les explications produites sur un modèle appris à ce que le modèle physique affirme. Le
chapitre~\ref{ch:xai} exploite directement cette possibilité.

Enfin, c'est le langage des interlocuteurs. Le BRGM développe et emploie sa propre famille de
modèles à réservoirs, dont GARDÉNIA~\cite{gardenia} pour la simulation des niveaux et des débits
d'un bassin et EROS~\cite{eros} pour les relations entre bassins, fédérés par une interface
Python commune. Un outil qui ignorerait cette pratique ne serait pas adopté, quelles que soient
ses performances.

\section{Les modèles à fonction de transfert et bruit}

L'approche retenue est celle des modèles \emph{transfer function-noise}, mise en \oe uvre par la
bibliothèque Pastas~\cite{pastas}. Son principe : le niveau simulé est la somme des
contributions de plusieurs facteurs de forçage, au minimum la recharge dérivée de la
précipitation et de l'évapotranspiration, chaque contribution étant obtenue par convolution du
forçage avec une \emph{fonction de réponse impulsionnelle} dont la forme et les paramètres sont
calibrés. Un modèle de bruit rend compte de la structure résiduelle.

L'intérêt pour la question posée est direct : la fonction de réponse est elle-même l'objet
d'intérêt hydrogéologique. Sa forme dit comment l'aquifère amortit et retarde le signal
climatique, son intégrale dit de combien le niveau varie pour une sollicitation donnée. Ce ne
sont pas des paramètres opaques, ce sont des grandeurs qui s'interprètent.

La plateforme donne accès à quatre modèles de recharge, six formes de réponse impulsionnelle,
deux modèles de bruit et deux solveurs, librement combinables.

\section{Calibration automatique et critères de qualité}

Choisir la bonne combinaison relève de l'expertise et de l'essai. Pour rendre cet essai
systématique, un mécanisme de calibration automatique construit une grille de configurations
candidates, calibre chacune, l'évalue et les classe.

L'évaluation ne repose pas sur un unique indicateur d'erreur mais sur les quatre critères
d'acceptation STOWA~\cite{stowa2012}, standard néerlandais d'évaluation des modèles à fonction
de transfert pour la gestion des eaux souterraines :

\begin{enumerate}
  \item la variance expliquée dépasse un seuil, fixé par défaut à \SI{70}{\percent} ;
  \item les résidus sont aléatoires, vérifié par un test de séquences ;
  \item le temps de réponse à \SI{95}{\percent} de la réponse indicielle reste inférieur à la
  moitié de la période de calibration ;
  \item le gain est statistiquement significatif, le rapport de sa valeur à son erreur type
  dépassant le seuil usuel.
\end{enumerate}

Le troisième critère mérite d'être souligné car il capture une erreur fréquente et discrète. Un
modèle peut afficher une excellente variance expliquée tout en ayant ajusté une dynamique lente
que la période d'observation ne permet pas de contraindre. Le critère l'écarte, là où un simple
indicateur d'erreur l'aurait accepté.

\section{L'apport hydrogéologique : des configurations par famille d'aquifère}

La grille de calibration n'explore pas l'espace des configurations à l'aveugle. Elle est amorcée
par des configurations recommandées selon la nature de l'aquifère observé, lue dans le
référentiel BDLISA rattaché à la station par l'entrepôt (chapitre~\ref{ch:entrepot}).

Cette table est le produit des échanges avec l'expertise hydrogéologique du BRGM. Elle traduit
en choix de modélisation ce qu'un spécialiste sait du comportement de chaque milieu.

\begin{center}
\begin{tabularx}{\textwidth}{@{}lllX@{}}
\toprule
\textbf{Famille} & \textbf{Réponse} & \textbf{Bruit} & \textbf{Justification hydrogéologique} \\
\midrule
Alluvial & Exponentielle & AR & Réponse rapide, temps de réponse court, système simple \\
Sédimentaire poreux & Gamma & AR & Réponse intermédiaire, sable et grès \\
Craie, double porosité & Gamma & AR & Forte inertie, temps de réponse long \\
Sédimentaire fracturé & Gamma & AR & Recharge complexe, modèle de recharge flexible \\
Karstique & Double exponentielle & ARMA & Deux temps de réponse : conduits et matrice \\
Socle et montagne & Gamma & AR & Fracturé, recharge non linéaire \\
Volcanique & Gamma & AR & Système hétérogène \\
Composite & Quatre paramètres & ARMA & Milieu hétérogène, modèle flexible \\
\bottomrule
\end{tabularx}
\end{center}

Le cas karstique est le plus parlant. Un système karstique se comporte comme deux réservoirs
couplés, des conduits qui réagissent en quelques heures et une matrice qui réagit en mois. Une
réponse impulsionnelle à un seul temps caractéristique ne peut pas représenter cette dualité,
d'où la double exponentielle et un modèle de bruit plus riche pour absorber ce que la structure
ne capte pas.

Cette table recoupe la typologie discutée lors de la réunion de cadrage avec les
hydrogéologues, qui distinguait les nappes inertielles à cycles longs, les nappes réactives à
effets de seuil et les nappes influencées par une rivière ou des prélèvements~\cite{reunion_aida}.

\begin{remarque}[title={Ce que cette table représente}]
Elle est un point de contact entre deux disciplines. Elle encode une connaissance qui ne
s'obtient pas par validation croisée et la rend disponible à un utilisateur qui ne la possède
pas. C'est une forme concrète de transfert d'expertise, et elle réduit le coût d'entrée à la
modélisation pour un non-spécialiste.
\end{remarque}

Ces configurations amorcent la recherche sans la contraindre : une grille de référence est
systématiquement évaluée en parallèle, de sorte qu'une station atypique n'est jamais enfermée
dans la recommandation de sa famille.

\section{Diagnostics et validation}

La calibration est accompagnée d'un dispositif de vérification qui va au-delà de l'ajustement.

\paragraph{Séparation calibration et validation.} Le modèle est calibré sur une période et
évalué sur une période distincte, avec les indicateurs usuels en hydrologie, efficience de
Nash-Sutcliffe, critère de Kling-Gupta et erreur quadratique moyenne, complétés par l'examen de
l'autocorrélation et de la normalité des résidus.

\paragraph{Diagnostics préalables.} Avant même la calibration, la chronique est examinée :
qualité et continuité des entrées, valeurs aberrantes, corrélations croisées avec le forçage,
décomposition du signal, analyse spectrale, analyse des tarissements. Ce faisceau permet de
détecter en amont les situations où la modélisation est vouée à l'échec, plutôt que de le
découvrir dans un mauvais ajustement.

\paragraph{Analyse multi-stations.} Les résidus peuvent être examinés conjointement sur
plusieurs stations, ce qui fait apparaître les structures communes, d'origine régionale ou
climatique, et les distingue des particularités locales.

\section{Scénarios prospectifs}

Le modèle calibré n'est pas seulement descriptif. Il permet de répondre à des questions de la
forme \emph{que se passerait-il si}, qui sont celles des gestionnaires.

Trois familles de scénarios sont outillées : l'ajout d'un prélèvement selon des profils
réalistes par type d'usage, alimentation en eau potable, irrigation ou industrie, chacun ayant
sa saisonnalité propre ; l'application d'une tendance climatique ; la modulation de l'intensité
d'un forçage existant.

Un mécanisme mérite d'être détaillé, celui des \textbf{bornes adaptatives}. Lorsqu'un
utilisateur simule un prélèvement, la question du débit plausible n'a pas de réponse
universelle : elle dépend de la transmissivité de l'aquifère, donc de la réponse calibrée. Le
rabattement attendu est estimé à partir de la réponse indicielle du modèle lui-même, et des
bornes physiquement plausibles en sont dérivées pour guider la saisie. L'utilisateur n'est ni
libre de simuler un pompage qui viderait la nappe en un mois, ni contraint par une borne
arbitraire identique partout.

\begin{acquisouvert}
\textbf{Acquis.} Un laboratoire de modélisation métier complet : calibration automatisée guidée
par l'expertise hydrogéologique, évaluation selon un standard reconnu, diagnostics préalables et
postérieurs, scénarisation prospective bornée par la physique du système. Il constitue la
référence à laquelle comparer toute approche apprise.

\textbf{Ouvert.} La comparaison systématique entre modèle métier et modèle appris, sur un large
échantillon de stations, reste à conduire. L'outillage des deux côtés est en place, la campagne
d'évaluation ne l'est pas. Un rapprochement avec les modèles à réservoirs du BRGM constitue une
seconde voie de comparaison, identifiée mais non explorée.
\end{acquisouvert}
